
les composantes irr\'eductibles de $X\cap Y^i=V(\varphi|Y^i)$ sont de dimension $\geq2$. Alors le morphisme $A^1_{\mathrm{num}}(Y)\longrightarrow A^1_{\mathrm{num}}(X)$ est injectif.}

En effet, comme on peut \'evidemment supposer $Y$ int\underline{\`egr}e, l'assertion r\'esulte de 7.4 (i), gra\^ce \`a la formule de projection.

\subsection*{BIBLIOGRAPHIE}
\begin{enumerate}[label={[\arabic*]}]
\item Abhyankar S. : Resolution of singularities of arithmetical surfaces, in : Arithmetical Algebraic Geometry, Purdue University, Harper-Row, New York 1965.
\item Mumford D : Pathologies III, Amer. J. Math 1967 p. 94-103 Amer. J. Math 1967 p.94-103.
\item Grothendieck A. . et Dieudonn\'e J. EGA, El\'ements de G\'eom\'etrie Alg\'ebrique, Pub. Math. Inst des Hautes Etudes Sci., 1960.
\item Grothendieck A. : FGA, Fondements de la g\'eom\'etrie alg\'ebrique, Extraits du S\'eminaire Bourbaki 1957-1962.
\item Grothendieck A. . et alii, SGA, S\'eminaire de g\'eom\'etrie alg\'ebrique, 1960-66, J.H.E.S.
\item Jussila O. : X, Formalisme des intersections sur les sch\'emas alg\'ebriques propres, Expos\'e X dans ce s\'eminaire.
\item Raynaud M. : XII, Un th\'eor\`eme de repr\'esentabilit\'e relative sur le foncteur de Picard, Expos\'e XII dans ce s\'eminaire.
\end{enumerate}



% Exposé XIV local notation
\providecommand{\Kb}{K^\bullet}
\providecommand{\Db}{D}
\providecommand{\Gtop}{\operatorname{Gr}_{\mathrm{top}}}
\providecommand{\Ftop}{\operatorname{Fil}_{\mathrm{top}}}


